\section{Interactive pool retrieval: protocol and security conditions}
\label{sec:interactive-pool-protocol}

This section specifies a replacement for \emph{pool retrieval}, conditional on
the surrounding Circuit OPRAM modules maintaining their existing invariants.
It is not a new proof of the complete ORAM construction.  The replacement uses
two-party secret-shared shuffling~\cite{cgp2020,osn2021}: the parties obtain
fresh XOR shares of a common random permutation of the records, while neither
party learns the composed permutation.  Both parties remain semi-honest, and
the application key and intermediate records remain secret shared.  The
ORAM architecture, recursive request processing, and eviction obligations
remain those of the surrounding construction~\cite{chan2017}.

\subsection{Interface and invariants}

Fix one recursive level and one batch.  Let $P$ be the public pool capacity,
$B$ the public batch size, and $u$ the number of UID bits at this level.  A
pool record contains a validity bit, UID, current leaf label, and payload.
The pool is a \emph{multiset stored in an array}; physical positions carry no
meaning needed by subsequent pool processing.  The following are required.

\begin{enumerate}
\item At most one valid record in the union of the pool and the lower tree
has a given UID.  Invalid records may contain arbitrary stale fields.  When
overlapping public paths are fetched, repeated appearances of one physical
tree slot are aliases, not additional independent live records.
\item Existing private conflict processing supplies a representative bit
$r_j$ for each request $j\in\{1,\ldots,B\}$.  There is at most one
representative per UID, and representatives have distinct UIDs.  Duplicate
group membership and result propagation remain secret.
\item Each requested live block is obtained from exactly one of the pool
and its fetched lower path.  Invalid and nonrepresentative requests are
masked.  The surrounding algorithm privately combines these possibilities
and clears each actual source once.
\item Later eviction and cleanup consume records through their secret
metadata, rather than interpreting a pool position as a fixed tree node.
If an implementation instead assigns persistent semantics to pool slots,
it must first change that interface or pay for inverse routing.
\end{enumerate}

These are interface conditions to verify in an implementation, not additional
leakage permissions.  In particular, the parties must not disclose whether
a request hit the pool, the lower tree, or a dummy.

\subsection{Match, tag, shuffle, and extract}

The tag alphabet is $\{0,1,\ldots,B\}$, with zero denoting an unselected row.
A tag has $t=\lceil\log_2(B+1)\rceil$ bits.  Tags identify public request
\emph{slots} in the current processing order; they do not disclose UIDs or
secret duplicate groups.  The following procedure is executed separately
at each recursive level, with fresh shuffling randomness.

\begin{quote}\small
\textbf{Input:} $P$ secret records $x_i$, and $B$ secret requests
$(q_j,r_j)$ in the surrounding circuit's current order.

\textbf{1. Private matching.} For every $i,j$, compute
\[
 e_{ij}=r_j\wedge\operatorname{valid}(x_i)
              \wedge[\operatorname{UID}(x_i)=q_j],\qquad
 h_j=\bigoplus_{i=1}^{P}e_{ij}.
\]
Set the tag of original row $i$ to the bitwise XOR of the public binary
constants $j$ for which $e_{ij}=1$.  There is at most one such $j$.

\textbf{2. Dummy completion.} Append one invalid record $d_j$ per request.
Give $d_j$ tag $j$ if $h_j=0$, and tag zero otherwise.  Its other fields
have the same fixed width and sharing representation as an ordinary record.
Nonrepresentatives have $h_j=0$ and therefore receive their own dummy.

\textbf{3. Secret shuffle.} Jointly shuffle the $n=P+B$ complete
$(\mathrm{tag},\mathrm{record})$ rows.  Each party samples a private uniform
permutation; the protocol applies their composition and refreshes shares.
Use a fresh protocol instance and fresh correlations for this invocation.

\textbf{4. Open tags only.} Reconstruct the shuffled tag column.  For each
$j$, both parties locally copy their record share at the unique row tagged
$j$ to output slot $j$.  Neither party reconstructs the record or its
validity bit.  Output shares remain secret.

\textbf{5. Retain the complement.} Both parties remove all $B$ nonzero-tag
rows locally and discard the public tags.  Keep the remaining $P$ records
in their current unordered arrangement.  Continue the existing private
path/pool result selection, duplicate-result routing, recursive updates,
and reinsertion/eviction schedule.
\end{quote}

At every invocation, each nonzero tag occurs exactly once: either its real
pool record or its dedicated dummy receives it.  Consequently the number
and pattern of opened tags do not reveal the number of hits.  A single
common dummy would be incorrect, as repeated opened positions would reveal
misses or request equality.  Tag assignment happens \emph{before} shuffling,
so no post-shuffle scan to rediscover dummy positions is needed.

\paragraph{Occupancy and clearing.}
Let $H$ be the number of representative pool hits.  Removal deletes $H$
original valid records and $B-H$ appended dummies.  The remaining $H$
appended dummies replace the removed real records, leaving exactly $P$
rows.  This extracts and clears the pool sources in one operation.  It is
not necessary to privately compact the survivors or return them to their
old slots.  Appending the surrounding algorithm's $B$ refreshed records
and masked dummy returns gives the usual $P+B$ candidate bound.  The
existing additional eviction workspace, candidate selection, and cleanup
are still required; this argument does not reduce their allowances.

\paragraph{Matching cost.}
Equality of two $u$-bit UIDs can be computed with $u-1$ ANDs, followed by
two ANDs for validity and representative gating.  Thus $BP(u+1)$ ANDs suffice
for this matching kernel.  The numerical ledger reserves one additional
AND per pair, charging $BP(u+2)$ as a conservative interface allowance.
Unique-hit aggregation is XOR.  Multiplication
of $e_{ij}$ by a public tag bit is either copying or zero, so forming tags
adds no payload-wide conditional selection.  This count excludes the
retained request-conflict module, path matching, and subsequent result
selection.  If valid-UID uniqueness is not established, XOR is not a
substitute for a general OR or a first-match selection circuit.

\subsection{Composition with recursion, paths, and updates}

Private request sorting remains part of the protocol.  The opened tag
returns a result to the same \emph{processing slot} that generated that tag.
It does not reveal the private sorting permutation.  Returning results to
the original query order still requires the existing private inverse
routing or an independently proved alternative.  Public tag extraction
cannot justify deleting those costs.

At a recursive level, several child requests may refer to one parent
position-map block.  Deduplication therefore uses that level's block UID,
not merely the original application key.  The parent block is fetched once;
its different child offsets are read and updated privately, updates are
coalesced, and one final parent block is committed.  This shuffle does not
make concurrent whole-block writes safe.  For a read batch, duplicate
queries can receive copies of the common result through existing secret
routing.  Supporting application writes additionally requires the chosen
operation order to be implemented privately within each duplicate group.

The pool and lower-tree results must be combined using secret hit/validity
flags.  A no-pool-hit dummy is a placeholder, not evidence of an absent
dictionary key.  Likewise, selecting a real pool row does not authorize
revealing that the lower path was unnecessary.  The existing fixed path
schedule is retained.  Nonrepresentatives use the surrounding scheme's
dummy-path rule, and collisions among public paths do not imply equality
of logical UIDs.  Shared physical slots on overlapping paths must have
their extraction/clear masks aggregated before a single state update.

Fresh leaf labels and updated metadata stay secret.  A removed valid pool
record is reintroduced only through the existing refreshed-return module;
an unselected record survives exactly once.  No later step may both retain
an old live copy and append a fresh live copy of the same UID.  These
conditions, together with correct coalescing of recursive updates, preserve
the valid-UID uniqueness needed by the next invocation.

\subsection{Security argument and its scope}

Consider a hybrid in which matching, conflict processing, share conversion,
and shuffling are replaced by their specified secure functionalities.
Conditional on the prior transcript, the pre-shuffle tag vector contains
exactly one copy of each $j\in\{1,\ldots,B\}$ and $P$ zeros.  A fresh
uniform permutation maps it to a uniform ordered placement of the $B$
distinct nonzero tags.  This distribution is independent of the hit vector,
the secret UIDs, and the original locations.  A simulator can therefore
sample the opened tag vector without learning any of those secrets.

One party knows its own private permutation, but the other party's
independent uniform permutation still hides the composition.  Secure
shuffling must refresh the shares: simply permuting recognizable active
labels or unchanged share strings is not this functionality.  Its output
record shares reveal no values when tags are opened.  Extraction and
complement removal are deterministic local functions of the public tags
and private shares, so they introduce no further openings.  Subsequent
secure modules consume these shares and reveal only outputs already
authorized by the surrounding protocol.  This is a conditional composition
argument for the replacement module, not a security or overflow proof for
the entire Circuit OPRAM template.

Every later batch must use a fresh hidden shuffle before disclosing further
positions.  Reusing a permutation permits cross-batch location linkage;
reusing its masks or correlations can expose differences of secret arrays.
Preprocessing can be performed ahead of time, but each correlation is
consumed once and must not be restored after a retry or rollback.  A retry
uses a new correlation.  In the semi-honest model, ``discard'' means that
the protocol does not reuse a correlation; privacy must hold even against
an adversary retaining its complete past view.

\subsection{Representations and explicit communication charges}

For the ordinary Boolean free-XOR backend, a bit $b$ is represented by a
garbler-held zero label $K^0$ and an evaluator-held active label
$K^b=K^0\oplus b\Delta$, where the permutation/color bit of $\Delta$ is
one.  Taking the color bits locally gives
\[
 [b]_G=\operatorname{col}(K^0),\qquad
 [b]_E=\operatorname{col}(K^b),\qquad
 [b]_G\oplus[b]_E=b.
\]
This adds no communication and discloses nothing outside each party's
existing view.  It does not itself promise independent fresh masks across
all output bits; the secure shuffle supplies the required refresh.  No
active label is sent to the garbler and no zero label is sent to the
evaluator during this conversion.

Returning an XOR-shared bit to an ordinary GC can be charged conservatively
as two party inputs followed by a free XOR: a 16-byte garbler input label,
two 16-byte OT message corrections for the evaluator input, and a 16-byte
per-OT matrix allowance for standard 128-bit IKNP extension.  This gives
\emph{64 bytes per logical re-input bit}, plus amortized base-OT setup and
transport padding.  This is an explicit accounting choice rather than an
optimal conversion claim.  Re-input every retained or extracted field that
the next GC actually consumes; the shuffle cost alone is not the complete
retrieval cost.  Persisting raw XOR shares instead of labels changes the
storage backend and requires charging subsequent touched-state re-input.

The color-bit argument applies to one-label-per-Boolean-bit encodings.
It does \emph{not} turn two opaque packed labels into shares of a 256-bit
value.  Accordingly, an ordinary Boolean model plus the explicit conversion
charge is the concrete baseline.  A two-label payload scenario remains
conditional on a specified packed encoding and secure conversion/update
interface.  One-hot and specialized garbling techniques may improve other
small-array modules~\cite{heath2024,heath2021}, but they do not remove the
representation or composition obligations here.

\paragraph{Shuffle cost.}
Let $R$ include all actual XOR-share bytes in a row, including its tag.
An $n$-input generalized Waksman network has
\[
 S(n)=n\lceil\log_2 n\rceil-2^{\lceil\log_2 n\rceil}+1
\]
switches.  In the construction of~\cite[Appendix A.3.1 and Figure 16]{osn2021},
a switch uses a 1-out-of-2 OT with two $2R$-byte candidate payloads.
Random-OT optimization reduces correction traffic from $4R$ to $2R$
bytes per switch.  Wire masks are chosen in advance, so these OTs can be
prepared in parallel; the data-dependent phase sends one masked input
array per private permutation.  Composing two private permutations gives
\[
 C_{\rm shuffle}=\underbrace{2nR}_{\rm online}
             +\underbrace{4S(n)R}_{\rm switch\ corrections}
             +C_{\rm OT}.
\]
Uniform permutations must be sampled first and then routed through the
network; independent fair switch bits do not sample uniform permutations.
All preprocessing communication belongs in the bandwidth/cost ledger even
when it lies outside the batch's critical path.  For $P=317$, $B=32$, and
$n=349$, $S(n)=2630$.  If a complete tagged row occupies 73 bytes, the first
two terms are 50,954 and 767,960 bytes per batch.  An illustrative 16-byte
OT-extension allowance per switch adds 84,160 bytes across the two
networks: 903,074 bytes per batch, or 28.22 decimal KB per request.  If the
73-byte width excludes the tag, rounding to 74 bytes gives 28.57 KB per
request instead.  Opening the tag column to both parties adds at most
$2\lceil nt/8\rceil$ bytes with bit packing: 524 bytes per batch here,
or 698 bytes if each tag is padded separately to one byte.  Metadata
matching, re-input, result routing, and the remaining ORAM work are
separate charges.
